\section{Security Deployment Context, Threat Model, and Ethics}

\subsection{Operational Threat Model}

The phishing URL detector in Question 2 can be interpreted as a first-stage classifier in a layered defense system. In this setting, adversaries attempt to maximize user trust while minimizing lexical anomalies that static rules can detect. Typical attacker strategies include:
\begin{itemize}
    \item Using HTTPS certificates to mimic legitimacy,
    \item Constructing long but semantically plausible paths,
    \item Using brand-like tokens with subtle lexical perturbations,
    \item Rotating domains rapidly to evade blacklist refresh cycles.
\end{itemize}

Given this threat model, a detector that optimizes only precision is insufficient, because missed phishing samples can have disproportionately high impact. The empirical results therefore favor models with stronger recall and balanced performance profiles.

\subsection{Deployment Pipeline Integration}

A practical deployment architecture for this assignment's model family can be structured into four stages:
\begin{enumerate}
    \item \textbf{Ingress filtering:} parse URL/email input and normalize formatting.
    \item \textbf{Feature extraction:} compute deterministic handcrafted features (or BoW vectors for email).
    \item \textbf{Primary classifier:} run Logistic Regression as balanced-risk scorer.
    \item \textbf{Escalation policy:} route uncertain or high-risk scores to secondary checks (sandboxing, DNS reputation, manual triage).
\end{enumerate}

This layered design reduces dependence on any single model assumption and supports iterative hardening as attacker behavior shifts.

\subsection{Decision-Policy Perspective}

In production systems, model outputs are not final decisions by themselves. They are inputs to policies that encode organizational risk appetite. A security-focused institution may prioritize recall and accept higher analyst workload from false positives. A high-volume consumer system may set a stricter threshold to preserve user experience, then compensate with auxiliary signals.

The report's ROC and PR curves provide the basis for selecting these operating thresholds. The current notebook uses default thresholds for comparability, but deployment should tune thresholds using cost-weighted validation.

\subsection{Monitoring and Drift Management}

Phishing and spam data are non-stationary. Model quality can decay even if training accuracy is strong. A complete operational plan should therefore monitor:
\begin{itemize}
    \item class prior drift,
    \item feature distribution drift (e.g., URL length and subdomain depth shifts),
    \item precision/recall degradation over time,
    \item error concentration by campaign or domain family.
\end{itemize}

Simple drift alarms can trigger retraining or temporary policy tightening. Without monitoring, even well-performing baselines become stale in adversarial environments.

\subsection{Ethical and Governance Considerations}

Although this project uses non-personalized lexical and structural features, deployment still raises governance concerns:
\begin{itemize}
    \item \textbf{False positives:} legitimate links or emails may be blocked, causing usability and trust issues.
    \item \textbf{Transparency:} users and administrators should receive understandable reason codes (e.g., suspicious token concentration).
    \item \textbf{Appeal and override:} systems should support safe review pathways for blocked content.
    \item \textbf{Security disclosure:} avoid exposing exact threshold logic publicly to reduce attacker adaptation.
\end{itemize}

A balanced design combines technical metrics with accountable policy controls, ensuring that detection quality does not come at the cost of opaque or brittle decisions.
